There are many ideas around (higher) category theory that are not always strictly formal.
These ideas are often only expressed in more or less precise English - or any other language with enough expressive power.
In other words, they only exist on the meta-level at this point in the discussion.
Of course, one would like to formalize these ideas to make them mathematically rigorous.
But this has not yet been achieved for all the important ideas presented in this section, albeit for some.
We split this section into two subsections:
\begin{enumerate}
\item[$\bullet$]
Internalization
\item[$\bullet$]
Higher Level Structure
\end{enumerate}
While the first subsection takes place in a $1$-categorical context, the second one is all about the higher categorical one.
As you will see, the idea of internalization also makes sense in a higher categorical context.
The reason for giving it its own section is rather its importance to us here.
We will concentrate on topoi and principal bundles there, alongside the core idea itself.
The section about higher-level structures is largely about building more intuition behind higher category theory and contains the important ideas of
\begin{enumerate}
\item[$\bullet$]
Categorification
\item[$\bullet$]
Oidification
\end{enumerate}
which we will then apply to a monoid in a subsubsection.
In this subsubsection, we will briefly discuss the further idea of
\begin{enumerate}
\item[$\bullet$]
Enrichment
\end{enumerate}
One last thing we want to explicitly remark on here is that it also contains the idea of operads, which are a nice tool to track the higher structure of higher categories and have their origins in homotopy theory.
